\section{Introduction}
\label{sec:intro}

The pursuit of Cryptographically Secure Software Obfuscation has evolved from foundational impossibility results to the vanguard of modern post-quantum mathematics. Originally postulated as theoretically impossible in the general Black-Box virtual \textit{Indistinguishability Obfuscation} (iO) paradigms by Barak et al., practical implementations shifted toward domain-specific \textit{White-Box Cryptography} (WBC) for digital rights management (DRM) and intellectual property shielding.

Early efforts in WBC, popularized by Chow et al. \cite{chow02}, relied heavily on Algebraic Encoding, Look-Up Table (LUT) network representations, and custom affine maskings to hide symmetric keys (like AES) directly integrated into the executable code. While theoretically formidable against strictly static analysis, these mechanisms ultimately succumbed to \textit{Differential Fault Analysis (DFA)}, \textit{BGE Attacks} \cite{bge04}, and dynamic runtime hooking vulnerabilities. 
If an adversary has infinite access to the CPU registers and can surgically modify memory streams mid-execution, purely algebraic LUTs provide zero cryptographic hardness beyond the obscurity of the topology itself.

Conversely, the discovery of \textit{Fully Homomorphic Encryption} (FHE) by Gentry \cite{gentry09} revolutionized cryptographic capabilities. FHE allows an entirely untrusted evaluator to process arbitrary mathematical circuits over encrypted data without possessing the decryption key. Modern derivatives relying on Ring Learning With Errors (RLWE) \cite{bv11} offer provable post-quantum security reductions to worst-case lattice problems. 

However, standard FHE specifically answers the \textbf{Honest-but-Curious} threat model (e.g., outsourced Cloud Computing verification). While the \textit{data vector} remains mathematically unbreakable, the \textit{control flow structure} and the literal \textit{algorithm} evaluating the data are fundamentally transparent to the server executing the homomorphic operations. FHE is not Obfuscation; it cannot hide a proprietary logical branching condition or mask a proprietary evaluating coefficient without invoking theoretically debilitating GGH-style Multilinear Maps \cite{ggh13}.

\subsection{Problem Statement}
This fundamental dichotomy leaves a critical void across modern cybersecurity applications: 
\begin{quote}
\textit{How can one attain robust, provable mathematical protection extending simultaneously across Data States and Control Flow Graph Logic in a Malicious White-Box Architecture, without incurring the physically infeasible computational overhead of generic Indistinguishability Obfuscation (iO)?}
\end{quote}

\subsection{Our Contribution: The Algebirc Hybrid Model}
In this paper, we address this void by introducing \textbf{Algebirc}, a novel hybrid cryptosystem proposing a structural bisection of the obfuscation methodology. Rather than forcing a singular primitive to secure orthogonal data abstractions, Algebirc utilizes the specific strengths of two disjoint cryptographic branches.

Our main contributions include:
\begin{itemize}
    \item \textbf{Polynomial RLWE Data Preservation:} We implement an optimized RLWE encryption scheme enveloping the memory architecture. We document the \textit{Multiplication Wall} flaw inherent in un-bootstrapped FHE within finite hardware matrices and bypass it utilizing a deterministic Key-Dependent \textit{Decrypt-Compute-Reencrypt (DCR)} primitive. While we strictly categorize DCR as a transitional Grey-Box mitigation rather than a mathematically perfect White-Box solution, it preserves static state indistinguishability while bypassing dimension-tensoring overheads.
    \item \textbf{Geometric Isogeny Control Mappings:} To protect the logical pathways natively susceptible to instruction-hook tracers during White-Box examination, Algebirc projects conditional boolean branches mapping into disjoint random walks over Genus-2 Hyperelliptic Curve Isogeny Graphs. This mechanism enables strict \textit{Oblivious Branchless Execution}.
    \item \textbf{Transparent Threat Contextualization:} We systematically analyze the strict limitations of substituting Black-Box principles into a White-Box model, providing exhaustive context defining the required \textit{Key-Dependent Message} (KDM-Security) \cite{bh20} properties preventing complete evaluation subversion across the DCR cycle.
\end{itemize}

\subsection{Organization}
This paper organizes conceptually as follows. Section \ref{sec:preliminaries} provides the dense theoretical foundations of RLWE hardness bounds alongside Jacobian Arithmetic for hyperelliptic geometries. Section \ref{sec:threat_model} rigorously maps the differentiation between the Black-Box Cloud and the target White-Box environments. Section \ref{sec:architecture} synthesizes the core algorithms spanning from generic initialization structures through the complete Hybrid routing. We present an exhaustive security evaluation over static analysis and Jacobian DLog resistance in Section \ref{sec:security}, before mapping tangible performance and hardware constraints in Section \ref{sec:performance}. Section \ref{sec:conclusion} concludes while analyzing the necessary transition trajectory toward TFHE Programmable Bootstrapping (PBS) architectures to accomplish total dynamic masking.
